\section{Mechanism evidence and historical interpretation}
\label{sec:mechanisms-history}

\subsection{A conditional finite-action mechanism}
% EXP-025/026; existing finite-action margin theory and completed diagnostics.
The synthetic and perturbation studies illustrate why absolute error and
decision error need not move together. Adding the same offset to every
predicted action value can change absolute fidelity but leaves pairwise
ordering unchanged. Relative perturbations can instead reverse an ordering
when the true alternatives are close. These are conditional finite-action
statements; they do not certify an available predictor's error bound.

Supplement~\ref{sec:explanatory-propositions} gives four explanatory
statements and proofs: common-mode invariance, strict STOP sign stability,
unique-optimum margin sufficiency, and margin-dependent regret with its
expectation bound. They require the same finite feasible action set for
prediction and evaluation. A bound uniform over actions at a state is
stronger than a small mean error across evaluation records. The proofs
therefore explain possible prediction--decision separation without claiming
that the fitted models satisfy their sufficient conditions.

The synthetic factorial and common-mode/relative-error perturbations
support this bounded interpretation (EXP-025 and EXP-026). The perturbations
reuse historical outputs from a learner whose features included realized
runtime. They therefore do not establish how a prospective learner behaves.
Their numerical mechanism grids are not part of the current prospective
comparisons.

\subsection{The margin explanation did not replicate}
The earlier FIT margin-bin association was not replicated by the frozen
validation40 VH3 comparison. That contrary result limits the empirical
mechanism claim: it would be misleading to present the favorable FIT trend
as an independently confirmed explanation. The present main argument uses
the prospective fidelity and external utility contrasts instead. Historical
D0--D3 decompositions likewise separate learned state selection from learned
action selection only by substituting oracle outcomes, and remain diagnostic
rather than deployable policies.

\subsection{What historical stratifications diagnose}
\label{sec:historical-diagnostic-dimensions}
Prediction and repair can be examined across training sizes, actions, image
groups, classes, costs, observations and error mechanisms. Each view asks
a different question about an aggregate result. The historical analyses
motivate these diagnostics, while their different learners and information
boundaries limit their relevance to the prospective comparisons.

\paragraph{Training scale.}
A scale comparison asks whether a pattern persists across matched amounts
of training information. The historical analyses compared discrete
training sizes and did not establish a scaling law.
Table~\ref{tab:f0f1} gives the registered remaining-action accuracy
contrasts across four scales. These contrasts concern the matched F0/F1
models, not an all-model learning curve or scale-dependent utility; their
exact evaluable-state and pair denominators remain unavailable.

\paragraph{Action composition.}
Per-action diagnostics ask whether an aggregate fidelity contrast is driven
by easy STOP predictions or by particular repairs. The historical R0/R1
comparison addressed this question for its own model and feature definitions.
Table~\ref{tab:coco-opportunity} reports the external
action-pair composition and oracle headroom. Those counts identify the
population behind a pooled metric; they are not per-action error or gain
estimates. A historical per-action ordering does not establish the ordering
for the prospective R1\_P model.

\paragraph{Image-group folds.}
Fold comparisons ask whether a pooled improvement depends on one held-out
image group partition. Historical fold-wise results used a learner and
feature boundary that differ from the prospective phase.
The current method retains grouped OOF fitting and paired image-group
resampling. Folds and seeds describe dependence and fitting variation;
they do not become additional independent test images. A fold configuration
specifies the evaluation design without establishing fold-wise performance.

\paragraph{Classes and population composition.}
Class diagnostics ask whether the aggregate conceals category-specific
reversals, with differing class denominators. The historical category
comparison does not establish an all-class advantage under the current
representation controls.
The prospective image--class construction and FIT-only target-rich analyses
instead make population composition explicit. Neither target-rich averages
nor the already selected examples in Appendix~\ref{sec:recovered-qualitative}
provide a new per-class performance curve or external category-wise test.

\paragraph{Costs and sparse budgets.}
Budget and cost diagnostics ask whether a controller locates beneficial
states and whether their gains repay intervention costs. D0--D3 further
separate state allocation from action choice
(Table~\ref{tab:selector_action_framework}). Historical FIT analyses used
this oracle decomposition and examined task-adapted responses to costs.
Table~\ref{tab:voc-r3-sparse} reports the six zero-cost R3 development
cells, and Table~\ref{tab:coco-sparse} gives both semantics for all six
external models. These results do not cover the full historical nonzero-cost
grid or support a claim of failure at every cost. V6 and COCO retain their
distinct utility criteria.

\paragraph{Observation.}
An already acquired probe changes the information available for predicting
the remaining actions. Historical probe analyses considered both accuracy
and utility under their own protocol. The prospective F0/F1 comparison
answers only the registered remaining-action accuracy question. Its eight
practical ties do not inherit the old utility gates or establish the value
of paying for A3. Other observation types and adaptive probes remain
outside this comparison.

\paragraph{Mechanism.}
Common-mode and relative perturbations isolate mathematical ways that
fidelity and decisions can separate. EXP-025/026 and the propositions have
this explanatory role. The empirical margin-stratification proposal also
needed replication, which the validation40 VH3 comparison did not provide.
The algebra remains valid while this proposed empirical explanation remains
unconfirmed. The historical validation result does not supply a new
prospective or external mechanism test.

\subsection{Earlier representation and objective routes}
The project executed early gain/post-action-loss models, sparse opportunity
decomposition, training-scale studies and frozen visual-feature routes.
The visual route included global and candidate-region features, ablations
and additional-seed/class-weighting checks (EXP-011--013). These were
executed historical experiments; the current prospective external panel
does not include those visual-feature models.

The original dense headroom gate failed, and subsequent sparse studies
addressed a separately frozen question. Early Q/SOAD, residual-to-decision,
R3 and A3-probe results retain their failed, mixed or negative gates.
Favorable isolated cells do not replace those terminals. The corrected
atom-bank accounting survives as a distinct mathematical and integrity
result; the invalid predecessor is not pooled into it. Historical
RETRO\_V4 learned results remain limited by future-action runtime and
other descriptor/geometry boundaries addressed by the prospective study.

\subsection{Panel use and procedural incidents}
Validation40 was used in early baseline/validation work before the later
locked diagnostic transaction. The latter's one-shot status applies to
that transaction, not to the lifetime of the image identity. Earlier
official-VOC-related metadata exposure, pilot-test metadata exposure and
replacement-panel assessments also limit claims of pristine confirmation.
The calibration panel likewise retains its historical role rather than
providing a fresh test population.

The public-data scientific comparisons here are the prospective VOC
development study and frozen COCO transfer. UAV work remains a data audit
without an external model result in this manuscript. Track B's planned R3-LCB and
benefit-value candidates were \texttt{NOT\_EVALUATED}: required
authenticated FIT dependencies were unavailable, and their \texttt{NO\_GO}
is not a scientific failure. The intermediate reviewer runtime-free plan
(EXP-032) remains execution-unverified and adds no assumed training batch.

Phase-qualified model-family and historical run totals describe the
project's scope. They do not count independent architectures, provide new
checkpoint certifications or establish scientific contributions.
